%=============================================================================!
% 9.3  REVERSAL AND VOLUME                                                    !
%=============================================================================!
\section{Reversal and volume}
\label{sec:in-reversal}

The exact motion can be run backwards (\cref{sec:ha-reversal}) and keeps
the volume of every region of phase space (\cref{sec:ha-liouville}). This
section shows that velocity Verlet keeps both properties exactly, for any
step, which the forward Euler method does not.

\subsection*{Running a step backwards}

We take a step from $(x, v)$ to $(x', v')$ by \cref{eq:in-kdk}, with
$v_{\frac12}$ the half-step velocity and $a$, $a'$ the accelerations at
$x$ and $x'$, then reverse the velocity and take one more step from $(x',
-v')$:
\begin{align*}
   -v' + \tfrac12\dt\,a' &= -\bigl(v' - \tfrac12\dt\,a'\bigr) = -v_{\frac12}
   && \text{the half kick undoes the last half kick,}\\
   x' + \dt\,(-v_{\frac12}) &= x' - (x' - x) = x
   && \text{the drift runs back to $x$,}\\
   -v_{\frac12} + \tfrac12\dt\,a &= -\bigl(v_{\frac12} - \tfrac12\dt\,a\bigr)
   = -v
   && \text{the last half kick undoes the first,}
\end{align*}
using $v' = v_{\frac12} + \frac12\dt\,a'$, $x' = x + \dt\,v_{\frac12}$ and
$v_{\frac12} = v + \frac12\dt\,a$. The step lands on $(x, -v)$: the
starting state with its velocity reversed, as for the exact motion. Since
every step can be undone this way, a whole run of velocity Verlet,
reversed, retraces itself; the tests of \texttt{mdlab} run 200 steps of
three atoms in nonlinear wells forwards and back and recover the start to
within $10^{-15}$, the rounding of the arithmetic. A forward Euler step does not
undo itself (\cref{ex:in-euler-reverse}).

\subsection*{Keeping volume}

For one coordinate, the half kick $(x, v) \to (x, v + \frac12\dt\,a(x))$
changes the velocity by an amount that depends only on the position, and
the drift $(x, v) \to (x + \dt\,v, v)$ changes the position by an amount
that depends only on the velocity: each is a shear, with Jacobian
determinant 1 (\cref{sec:ha-symplectic}). Velocity Verlet is three such
shears in turn, so it keeps area in the phase plane, as \cref{fig:in-shadow}
will show: its states for the spring lie on a closed curve, where those of
forward Euler spiral outwards (\cref{fig:ha-maps}), the test of
\cref{sec:os-phase}.

For $N$ atoms the same holds in the $6N$ dimensions of phase space. In a
kick, the new momenta are the old plus amounts that depend only on the
positions, and the positions do not change; its Jacobian matrix, with the
positions listed before the momenta, has 1 on its diagonal and zeros
everywhere above the diagonal. A square matrix with only zeros on one side
of its diagonal has as determinant the product of its diagonal entries, a
property of determinants we take on trust, here 1. The drift is the same
with the roles exchanged, so every step of velocity Verlet keeps volume in
phase space, as the exact flow does (\cref{sec:ha-liouville}).

\subsection*{The full condition}

With more than one coordinate, keeping volume is not all that the exact
flow keeps, and saying what else it keeps needs products of matrices.

\begin{toolbox}[Products of matrices]
The product $\vb{A}\vb{B}$ of two $n\times n$ matrices is the matrix whose
column $b$ is $\vb{A}$ times column $b$ of $\vb{B}$, so that its entry in
row $a$ and column $b$ is $(\vb{A}\vb{B})_{ab} = \sum_c A_{ac}B_{cb}$, the
entries of row $a$ of $\vb{A}$ times those of column $b$ of $\vb{B}$,
added. For example,
\begin{equation*}
   \begin{pmatrix} 1 & 2\\ 0 & 1 \end{pmatrix}
   \begin{pmatrix} 1 & 0\\ 3 & 1 \end{pmatrix}
   = \begin{pmatrix} 7 & 2\\ 3 & 1 \end{pmatrix} ,
   \qquad
   \begin{pmatrix} 1 & 0\\ 3 & 1 \end{pmatrix}
   \begin{pmatrix} 1 & 2\\ 0 & 1 \end{pmatrix}
   = \begin{pmatrix} 1 & 2\\ 3 & 7 \end{pmatrix} ,
\end{equation*}
so the order of a product matters. Since each column of $\vb{A}\vb{B}$ is
$\vb{A}$ applied to a column of $\vb{B}$, $(\vb{A}\vb{B})\vb{u} =
\vb{A}(\vb{B}\vb{u})$ for every vector $\vb{u}$: the product is the map
$\vb{B}$ followed by the map $\vb{A}$. The transpose $\vb{A}^{\mathsf T}$
is $\vb{A}$ with rows and columns exchanged, $(\vb{A}^{\mathsf T})_{ab} =
A_{ba}$, and exchanging the rows and columns of a product reverses its
order, since
\begin{equation*}
   \bigl((\vb{A}\vb{B})^{\mathsf T}\bigr)_{ab} = (\vb{A}\vb{B})_{ba}
   = \sum_c A_{bc}B_{ca} = \sum_c(\vb{B}^{\mathsf T})_{ac}
   (\vb{A}^{\mathsf T})_{cb} = (\vb{B}^{\mathsf T}\vb{A}^{\mathsf T})_{ab} .
\end{equation*}
A large matrix may be written in blocks, smaller matrices set side by side
and one above another. The sum $\sum_c A_{ac}B_{cb}$ then splits into sums
over the columns of each block, and two such matrices multiply block by
block as $2\times2$ matrices of numbers do, keeping the order of every
product of blocks. The trace of a square matrix is the sum of its diagonal
entries; for the $2\times2$ matrix of \cref{sec:os-modes}, with $a$, $b$
in the first row and $c$, $d$ in the second, the characteristic equation
$(a - \lambda)(d - \lambda) - bc = 0$ multiplies out to $\lambda^2 - (a +
d)\lambda + (ad - bc) = 0$, whose coefficients are the trace and the
determinant.
\end{toolbox}

We write a point of phase space as $\Gamma = (\vb{q}, \vb{p})$ and the
Jacobian matrix of a step map as $\vb{J}$, so that, by the first-order
rule of \cref{sec:en-gradient}, a small change $\delta\Gamma$ of the
starting point becomes the change $\vb{J}\,\delta\Gamma$ after the step.
The flow of Hamilton's equations over any time, and so any map that is
the flow of some Hamiltonian, satisfies (a result we take on trust)
\begin{equation}
   \vb{J}^{\mathsf T}\boldsymbol{\Omega}\,\vb{J} = \boldsymbol{\Omega} ,
   \qquad
   \boldsymbol{\Omega} = \begin{pmatrix} \vb{0} & \vb{I}\\ -\vb{I} & \vb{0}
   \end{pmatrix} ,
   \label{eq:in-symplectic}
\end{equation}
where $\vb{I}$ and $\vb{0}$ are the identity and zero matrices of the size
of $\vb{q}$, placed as four blocks. A map that satisfies
\cref{eq:in-symplectic} is called symplectic\cite{LeimkuhlerMatthews2015}.
For one coordinate the condition reduces to $\det\vb{J} = 1$
(\cref{ex:in-omega}), which is why \cref{sec:ha-symplectic} could use
keeping area as its definition; with more coordinates it implies keeping
volume but is stronger. Half kicks and drifts are flows, over a time
$\frac12\dt$ or $\dt$, of the Hamiltonians $U(\vb{q})$ and $K(\vb{p})$
alone, so they satisfy it. So
does any sequence of them: two maps in turn take $\delta\Gamma$ to
$\vb{J}_1\delta\Gamma$ and then to $\vb{J}_2\vb{J}_1\delta\Gamma$, so the
Jacobian matrix of the pair is the product $\vb{J}_2\vb{J}_1$, and
\begin{equation*}
   (\vb{J}_2\vb{J}_1)^{\mathsf T}\boldsymbol{\Omega}\,\vb{J}_2\vb{J}_1
   = \vb{J}_1^{\mathsf T}\bigl(\vb{J}_2^{\mathsf T}\boldsymbol{\Omega}
   \vb{J}_2\bigr)\vb{J}_1 = \vb{J}_1^{\mathsf T}\boldsymbol{\Omega}\vb{J}_1
   = \boldsymbol{\Omega} .
\end{equation*}
Velocity Verlet, a half kick, a drift and a half kick, is therefore
symplectic.

\begin{keyidea}
Reversing the velocities after a step of velocity Verlet and stepping
again returns the start, with its velocities reversed, without error, and
every step, being kicks and a drift,
keeps volume in phase space and satisfies the full symplectic condition
$\vb{J}^{\mathsf T}\boldsymbol{\Omega}\vb{J} = \boldsymbol{\Omega}$.
\end{keyidea}

\begin{selfcheck}
For a single coordinate, write the Jacobian matrix of the half kick
$(x, v) \to (x, v + \frac12\dt\,a(x))$ and find its determinant.
\end{selfcheck}
